%%%%%%%%%%%%%%%%
%%% lev 2 %%%
%%%%%%%%%%%%%%%%

\Chapter{2}

\Verse{1}
{
	\hP{וְנֶפֶשׁ כִּי תַקְרִיב קרְבַּן מִנְחָה לַיהֹוָה סֹלֶת יִהְיֶה קרְבָּנוֹ וְיָצַק עָלֶיהָ שֶׁמֶן וְנָתַן עָלֶיהָ לְבֹנָה׃}
}
{
	\eP{
		“And a person who will make an offering of a grain
		offering to YHWH, his offering shall be of fine flour, and
		he shall pour oil on it and put frankincense on it.
	}
}
{
	\aB{
		When you people make bread present to YHWH, use best flour.
		Pour oil, put frankincense. Very simple.

		Don't eat the frankincense. Frankincense not food. Frankincense is air freshener.

		Good smell for YHWH.
	}
}

\Verse{2}
{
	\hP{וֶהֱבִיאָהּ אֶל בְּנֵי אַהֲרֹן הַכֹּהֲנִים וְקָמַץ מִשָּׁם מְלֹא קֻמְצוֹ מִסּלְתָּהּ וּמִשַּׁמְנָהּ עַל כּל לְבֹנָתָהּ וְהִקְטִיר הַכֹּהֵן אֶת אַזְכָּרָתָהּ הַמִּזְבֵּחָה אִשֵּׁה רֵיחַ נִיחֹחַ לַיהֹוָה׃}
}
{
	\eP{
		And he shall bring it to the sons of Aaron, the priests.
		And he shall take the fill of his fist from there: from
		its fine flour and from its oil in addition to all of its
		frankincense. And the priest shall burn a representative
		portion of it to smoke at the altar, an offering by fire
		of a pleasant smell to YHWH.
	}
}
{
	\aB{
		Bring to Aaron son.

		Priest take handful, burn on altar.

		Burn all frankincense.\footnote{Don't eat the frankincense.}

		Good smell for YHWH.
	}
}

\Verse{3}
{
	\hP{וְהַנּוֹתֶרֶת מִן הַמִּנְחָה לְאַהֲרֹן וּלְבָנָיו קֹדֶשׁ קדָשִׁים מֵאִשֵּׁי יְהֹוָה׃}
}
{
	\eP{
		And the remainder from the grain offering is Aaron’s and
		his sons’, the holy of holies from YHWH’s offerings by
		fire.
	}
}
{
	\aB{
		Leftover is for Aaron son.

		Olive oil flatbread is holy of holy leftover.
	}
}

\Verse{4}
{
	\hP{וְכִי תַקְרִב קרְבַּן מִנְחָה מַאֲפֵה תַנּוּר סֹלֶת חַלּוֹת מַצֹּת בְּלוּלֹת בַּשֶּׁמֶן וּרְקִיקֵי מַצּוֹת מְשֻׁחִים בַּשָּׁמֶן׃}
}
{
	\eP{
		“And if he will make an offering of an oven-baked grain
		offering: fine flour, unleavened cakes mixed with oil and
		unleavened wafers with oil poured on them.
	}
}
{
	\aB{
		When you make baked good, use good flour.
	
		Olive oil flatbread and olive oil cracker.

		Text never say olive oil.

		Any oil fine: olive oil, canola oil, vegetable oil, soybean oil,
		corn oil, sunflower oil, avocado oil, sesame oil, coconut oil,
		rapeseed oil, grapeseed oil, all good.

		No peanut oil. Priest maybe have allergy. Don't kill priest.

		No palm oil. Too industrial.

		No butter, ghee, lard. Those not oil.

		Anyone who bring lard cracker to YHWH is getting cut off from his people.

		\heb{היה}.
	}
}

\Verse{5}
{
	\hP{וְאִם מִנְחָה עַל הַמַּחֲבַת קרְבָּנֶךָ סֹלֶת בְּלוּלָה בַשֶּׁמֶן מַצָּה תִהְיֶה׃}
}
{
	\eP{
		“And if your offering is a grain offering on a griddle, it
		shall be fine flour, mixed with oil, unleavened.
	}
}
{
	\aB{
		Grain offering on griddle, very simple.

		Good flour. Olive oil. Mix together.
	
		Don't put yeast on offering, people.
	
		\heb{היה}.
	}
}

\Verse{6}
{
	\hP{פָּתוֹת אֹתָהּ פִּתִּים וְיָצַקְתָּ עָלֶיהָ שָׁמֶן מִנְחָה הִוא׃}
}
{
	\eP{
		Break it into pieces and pour oil on it. It is a grain
		offering.
	}
}
{
	\aB{
		Break into pieces and drizzle more oil.

		It is grain offering.
	}
}

\Verse{7}
{
	\hP{וְאִם מִנְחַת מַרְחֶשֶׁת קרְבָּנֶךָ סֹלֶת בַּשֶּׁמֶן תֵּעָשֶׂה׃}
}
{
	\eP{
		“And if your offering is a grain offering from a pan, it
		shall be made of fine flour in oil.
	}
}
{
	\aB{
		% The "Looks like p, is secretly w" Family
		% Old English Wynn: ƿ (perfect choice)
		%
		% The "Looks like a, is secretly o/ɔ" Family
		% Syriac waw: ܘ
		% IPA Open back unrounded vowel: ɑ (best)
		%
		% The "Looks like n, is secretly k" Family
		% Old North Arabian: 𐪋
		% Old South Arabian: 𐩫
		% Ge'ez descendant: ከ (best)
		If your offering cooked in ƿɑከ, you have wisdom.

		Always use ƿɑከ for grain offering.

		Especially fried grain%
			\footnote{Lit. \chineseB{\ruby{米}{mai5}} in the original,
			presumably a rendering of the term ``maize,'' a cereal grain
			like wheat, barley, oats, rye, and millet.
		} offering.

		Only use ƿɑከ for food, pan%
			\footnote{
				Literally \ruby{\chineseB{炒}}{\heb{צאו}} \ruby{\chineseB{飯}}{\heb{פאנ}} in the original,
				a complex arrangement of glyphs with annotations suggesting a long history of
				scribal confusion concerning the meaning of the expression. The Heb. annotation appears to
				say ``chow pan.'' We have rendered it as ``food, pan'' with the added comma and more
				standard culinary terminology to clarify the syntactic and semantic nature of text.
			} the%
			\footnote{Lit. \heb{ה} in the original.}
			only ƿɑከ for Him.%
			\footnote{Lit. ``with He'' in the original. Meaning of Eng. uncertain.}

		Use oil in the\footnote{Lit. \heb{ה} in the original.}
		ƿɑከ, make ƿɑከ have Him.\footnote{Lit. ``have He'' in the original. Meaning of Eng. uncertain.}

		You%
		\footnote{Lit. \ruby{\chineseB{曰}}{yuē} in the original text.}
		listen, this not
		joke,%
		\footnote{Lit. \ruby{\chineseB{曰}}{joek5} in the original text.}
		this Word%
		\footnote{%
			Lit. \ruby{\chineseB{曰}}{\ruby{lou}{\chineseB{姥}}\ruby{gou}{\chineseB{糕}}} in the original text,
			the annotations appear to preserve a scribal tradition of interpreting this as the Grk. λογος.
		}
		of God.%
		\footnote{Lit. \ruby{\chineseB{曰}}{/*ɢʷad/} in the original text.}
	}
}

\Verse{8}
{
	\hP{וְהֵבֵאתָ אֶת הַמִּנְחָה אֲשֶׁר יֵעָשֶׂה מֵאֵלֶּה לַיהֹוָה וְהִקְרִיבָהּ אֶל הַכֹּהֵן וְהִגִּישָׁהּ אֶל הַמִּזְבֵּחַ׃}
}
{
	\eP{
		“And you shall bring the grain offering that one will make
		from these things to YHWH, and one shall bring it forward
		to the priest, and he shall bring it over to the altar.
	}
}
{
	\aB{
		Bring to YHWH.

		That means give to priest. Priest bring to altar.

		Don't just walk up to altar yourself. This is Tent of Meeting. \heb{היה}.
	}
}

\Verse{9}
{
	\hP{וְהֵרִים הַכֹּהֵן מִן הַמִּנְחָה אֶת אַזְכָּרָתָהּ וְהִקְטִיר הַמִּזְבֵּחָה אִשֵּׁה רֵיחַ נִיחֹחַ לַיהֹוָה׃}
}
{
	\eP{
		And the priest shall lift from the grain offering a
		representative portion of it and burn it to smoke at the
		altar, an offering by fire of a pleasant smell to YHWH.
	}
}
{
	\aB{
		Priest take some and burn on altar.

		Good smell for YHWH.
	}
}

\Verse{10}
{
	\hP{וְהַנּוֹתֶרֶת מִן הַמִּנְחָה לְאַהֲרֹן וּלְבָנָיו קֹדֶשׁ קדָשִׁים מֵאִשֵּׁי יְהֹוָה׃}
}
{
	\eP{
		And the remainder from the grain offering is Aaron’s and
		his sons’, the holy of holies from YHWH’s offerings by
		fire.
	}
}
{
	\aB{
		Leftover for Aaron son.

		Holy of holy from YHWH.
	
		\heb{יהי}.
	}
}

\Verse{11}
{
	\hP{כּל הַמִּנְחָה אֲשֶׁר תַּקְרִיבוּ לַיהֹוָה לֹא תֵעָשֶׂה חָמֵץ כִּי כל שְׂאֹר וְכל דְּבַשׁ לֹא תַקְטִירוּ מִמֶּנּוּ אִשֶּׁה לַיהֹוָה׃}
}
{
	\eP{
		Every grain offering that you will bring forward to YHWH
		shall not be made with leavening, because all leaven and
		all honey: you shall not burn any of it to smoke as an
		offering by fire to YHWH.
	}
}
{
	\aB{
		No yeast.

		No honey on altar.

		Who put honey on altar?

		\heb{היה}.

%	\footnote{
%		Every grain offering … shall not be made with leavening. That is, no
%		grain offering shall be made with leaven. See the comment on Exod
%		Footnote 29:2.
%	}
	}
}

\Verse{12}
{
	\hP{קרְבַּן רֵאשִׁית תַּקְרִיבוּ אֹתָם לַיהֹוָה וְאֶל הַמִּזְבֵּחַ לֹא יַעֲלוּ לְרֵיחַ נִיחֹחַ׃}
}
{
	\eP{
		You shall bring them forward to YHWH as an offering of a
		first thing, but they shall not go up to the altar as a
		pleasant smell.
	}
}
{
	\aB{
		Bring to YHWH.

		First thing of harvest.

		Don't burn it.

		This not for smell.
	}
}

\Verse{13}
{
	\hP{וְכל קרְבַּן מִנְחָתְךָ בַּמֶּלַח תִּמְלָח וְלֹא תַשְׁבִּית מֶלַח בְּרִית אֱלֹהֶיךָ מֵעַל מִנְחָתֶךָ עַל כּל קרְבָּנְךָ תַּקְרִיב מֶלַח׃}
}
{
	\eP{
		And you shall sprinkle every offering of a grain offering
		with salt. And you shall not let the salt of your God’s
		covenant cease from on your grain offering; you shall
		bring salt on all your offerings.
	}
}
{
	\aB{
		Put salt on offering. Food need seasoning.

		Covenant say salt. YHWH know seasoning.

		Imagine YHWH make universe and you bring him bland cracker. \heb{היה}.

%	\footnote{
%		sprinkle every offering of a grain offering with salt. The requirement
%		to use salt with these offerings is now memorialized by the custom of
%		putting salt on bread when one says the blessing over bread (hammî’) at
%		the beginning of a meal.
%	}
	}
}

\Verse{14}
{
	\hP{וְאִם תַּקְרִיב מִנְחַת בִּכּוּרִים לַיהֹוָה אָבִיב קָלוּי בָּאֵשׁ גֶּרֶשׂ כַּרְמֶל תַּקְרִיב אֵת מִנְחַת בִּכּוּרֶיךָ׃}
}
{
	\eP{
		“And if you will bring a grain offering of firstfruits to
		YHWH, you shall bring the grain offering of your
		firstfruits ripe, parched with fire, groats of fresh
		grain.
	}
}
{
	\aB{
		First grain of harvest. Very fancy.

		Cook over fire. If grain still pale, what you doing?

		\heb{היה}.
	}
}

\Verse{15}
{
	\hP{וְנָתַתָּ עָלֶיהָ שֶׁמֶן וְשַׂמְתָּ עָלֶיהָ לְבֹנָה מִנְחָה הִוא׃}
}
{
	\eP{
		And you shall put oil on it and set frankincense on it. It
		is a grain offering.
	}
}
{
	\aB{
		Put oil on it.

		Then frankincense okay.

		It is grain offering.
	}
}

\Verse{16}
{
	\hP{וְהִקְטִיר הַכֹּהֵן אֶת אַזְכָּרָתָהּ מִגִּרְשָׂהּ וּמִשַּׁמְנָהּ עַל כּל לְבֹנָתָהּ אִשֶּׁה לַיהֹוָה׃}
}
{
	\eP{
		And the priest shall burn a representative portion of it
		to smoke, from its groats and from its oil, in addition to
		all of its frankincense, an offering by fire to YHWH.
	}
}
{
	\aB{
		Now priest burn some of fried grain offering.

		Some oil and all frankincense.

		Nobody eat incense.

		How many time I have to tell you.

		Who put incense in rice?

		For love of YHWH.\footnote{Lit. \chineseA{愛} \heb{יה} in the original text.}
	}
}
